% Original-source acquisition and scope comparison, 19 November 2026.
% Unnumbered replacement for the stale source-access paragraph in A.6.
\paragraph{The original 1990 and 1997 sources.}
The three original Crelle articles have now been opened from the
public G\"ottingen scans. Nakamura's volume~405 paper studies the
quotient with metabelian geometric kernel; Theorems~A--B retain
additional field or exceptional-unit hypotheses, while Theorem~C
recovers a multiplicative-difference invariant
\cite{A567_Nakamura1990Metabelian}.
His distinct volume~411 paper, Definition~3.3 and Theorem~3.4,
pp.~209--210, proves smooth cusp recognition over number fields using
self-normalizing procyclic subgroups, open arithmetic normalizer
image and the cyclotomic character. This confirms the original
provenance of the scoped smooth-cusp result used through Stix above.
Theorems~5.10 and~6.1, p.~215, reconstruct isomorphism of genus-zero
objects; Theorem~5.10 also realizes the induced cusp permutation.
They do not state that every prescribed group isomorphism is
induced by that curve isomorphism \cite{A567_Nakamura1990}.

Ihara--Nakamura's construction starts with a maximally degenerate
curve whose irreducible components are smooth projective lines.
Its canonical parametrization requires both a tangential structure
and an ordering of the nodes. For the chosen diagonal deformation,
Theorem~2$'$ compares kernels of Galois actions on tangential
pro-$\mathcal C$ fundamental groupoids.
Theorem~3B gives the universal pro-$\ell$ monodromy fixed-field
inclusion, and Corollary~4.2.2 gives equality for $n\geq1$ using
the previously known opposite inclusion
\cite[\S\S1.2--1.4, 3.6, 4.1--4.2, pp.~125--127,146--148]{A567_IharaNakamura1997}.
Immediately after Corollary~4.2.2, the article credits Takao with
the unpunctured case, without reproducing his proof.
Thus Oda's fixed-field problem is distinct from reconstruction in
049(b), as also explained by the opened Collas--Philip paper.
L6 identifies a relevant later programme; the tangential,
compactification and arithmetic-presentation assertions of
049(a--c) require the additional inputs and arguments supplied here.
